\documentclass[10pt,letterpaper]{article}
\usepackage[margin=.78in]{geometry}
\usepackage{amsmath,amssymb,amsthm}
\usepackage[hidelinks]{hyperref}

\newtheorem{theorem}{Theorem}
\newtheorem{proposition}{Proposition}

\title{Natural density of the prime--exponent divisibility sequence}
\author{Pioneer project}
\date{28 November 2026}

\begin{document}
\begin{center}
{\Large\bfseries Natural density of the prime--exponent divisibility sequence}\par\smallskip
Pioneer project \quad 28 November 2026
\end{center}

\begin{abstract}
We prove the density-zero conjecture for OEIS A392819, whose terms \(n\)
satisfy \(p+v_p(n)\mid n\) for every prime divisor \(p\) of \(n\).
A large prime divisor forces an additional divisibility by a modulus of
size at least \(p^2\). Integers with only small prime divisors contribute
zero density. We also give an odd term with two distinct prime divisors.
\end{abstract}

For each positive integer \(n\) and prime \(p\), let \(v_p(n)\) denote
the exponent of \(p\) in the prime factorization of \(n\). Set
\[
  S=\{n\ge 1 : p+v_p(n)\mid n \text{ for every prime }p\mid n\},
  \qquad S(X)=\#\{n\in S: n\le X\}.
\]
Felix Huber introduced the terms of \(S\) as OEIS A392819 and conjectured
that the set has natural density zero \cite{a392819}. The same set is
the fixed-point set of the least-common-multiple map in A393351
\cite{a393351}.

\begin{theorem}
\[
  \lim_{X\to\infty}\frac{S(X)}{X}=0.
\]
More precisely, for integers \(R\ge 2\), \(K\ge 1\), and \(X\ge 1\),
\[
 \frac{S(X)}{X}
 \le \frac{1}{X}\prod_{m=2}^{R-1}m^K
      +\frac{R}{2^K}+\frac{2}{R-1},
\]
with an empty product interpreted as one.
\end{theorem}

\begin{proof}
Suppose \(n\in S\) has a prime divisor \(p\ge R\). If \(v_p(n)\ge 2\),
then \(p^2\mid n\). If \(v_p(n)=1\), the defining condition gives
\(p+1\mid n\). Since \(p\) and \(p+1\) are coprime, this gives
\(p(p+1)\mid n\). The number of multiples of a positive integer \(d\)
up to \(X\) is at most \(X/d\). Counting by the two possible moduli
for each prime bounds the number of these terms by
\[
 X\sum_{p\ge R}\left(\frac1{p^2}+\frac1{p(p+1)}\right)
 \le 2X\sum_{m=R}^{\infty}\frac1{m^2}
 \le 2X\sum_{m=R}^{\infty}
   \left(\frac1{m-1}-\frac1m\right)
 =\frac{2X}{R-1}.
\]
The first sum is over primes. The last inequality uses
\(m(m-1)\le m^2\).

Every remaining term has all prime factors smaller than \(R\).
If any such prime occurs to exponent at least \(K\), the term belongs
to a class of multiples of \(p^K\). Summing their counts over
\(2\le p<R\) bounds the total by \(RX/2^K\). Otherwise each exponent
is smaller than \(K\), and the term is at most
\(\prod_{m=2}^{R-1}m^K\), since every prime factor is included in
this product with exponent \(K\). The claimed inequality follows.
First let \(X\to\infty\) with \(K,R\) fixed, then let \(K\to\infty\),
and finally let \(R\to\infty\). The lower bound \(S(X)/X\ge 0\)
gives the limit.
\end{proof}

\begin{proposition}
Every odd term of A392819 is a square. An odd term need not be a
prime power.
\end{proposition}

\begin{proof}
For each prime divisor \(p\) of an odd term \(n\), the integer
\(p+v_p(n)\) divides \(n\) and is therefore odd. As \(p\) is odd,
\(v_p(n)\) is even. This proves the square claim. For the second claim,
take \(5625=3^2\cdot5^4\). Both \(3+2=5\) and \(5+4=9\) divide
\(5625\), so it is an odd term with two prime divisors. This corrects
the assertion in the A392819 comments that every odd term is a prime
power \cite{a392819}.
\end{proof}

The accompanying verifier computes \(S(X)\) directly by trial division
and checks the factorization and the two divisibilities for \(5625\).
It gives \(S(100)=8\), \(S(1000)=30\), \(S(10000)=122\), and
\(S(50000)=313\).

\begin{thebibliography}{9}
\bibitem{a392819}
OEIS Foundation Inc., \emph{The On-Line Encyclopedia of Integer Sequences},
entry A392819, submitted by Felix Huber, 23 February 2026.
\url{https://oeis.org/A392819}.

\bibitem{a393351}
OEIS Foundation Inc., \emph{The On-Line Encyclopedia of Integer Sequences},
entry A393351. \url{https://oeis.org/A393351}.
\end{thebibliography}
\end{document}
